\section{Exact complexity of a single physical product}\label{sec:single-product}
A ratio bound can be small even when evaluating an individual factor's
exact envelope is difficult. For one factor the ratio is identically
one wherever its gap is positive. The next result concerns computing
that factor's envelope, not loss caused by separating interacting terms.
All complexity statements use rational binary encodings and bit time.

\begin{theorem}[Midpoint decision on narrow unequal boxes]
\label{thm:single-product-hardness}
Given rational $u_i>1$ and $q$, deciding whether
\[
 \vex_B m(\bar x)\le q,\qquad m(x)=\prod_{i=1}^nx_i,
 \quad B=\prod_i[1,u_i],\quad \bar x_i=(1+u_i)/2,
\]
is NP-complete. For every fixed rational $\eta>0$ it remains
NP-complete when $u_i\le1+\eta$ for all $i$.
\end{theorem}
\begin{proof}
Reduce from PARTITION, the classical NP-complete integer partition
problem \citep{Karp1972}. Let its positive integers be $a_1,\ldots,a_n$
with total $A$. Doubling them permits us to assume $A$ even and at least
two. Fix an integer $K\ge1$ and put
\[
 \epsilon=(16KA^3)^{-1},\qquad u_i=1+\epsilon a_i.
\]
These boxes satisfy $1<u_i\le1+1/(16KA^2)<2$.
In normalized endpoint variables $X_i=1+\epsilon a_iZ_i$, the midpoint
means are $\E Z_i=1/2$. Write $S(z)=\sum_i a_i z_i$. On a binary vertex,
\begin{equation}\label{eq:partition-expansion}
 \prod_i(1+\epsilon a_iz_i)
 =1+\epsilon S(z)+\frac{\epsilon^2}{2}
             \left(S(z)^2-\sum_i a_i^2z_i\right)+R(z).
\end{equation}
All remainder coefficients are nonnegative. Its order-$j$ elementary
symmetric sum is at most $A^j$, and $\epsilon A\le1/16$, so
\begin{equation}\label{eq:partition-remainder}
 0\le R(z)\le\frac{(\epsilon A)^3}{1-\epsilon A}
     \le2\epsilon^3A^3=\epsilon^2/(8K)\le\epsilon^2/8.
\end{equation}
Every admissible law has $\E S=A/2$, whence
\[
 \E m(X)=b+\frac{\epsilon^2}{2}\operatorname{Var}(S)+\E R(Z),
 \qquad b=1+\frac{\epsilon A}{2}
       +\epsilon^2\left(\frac{A^2}{8}-\frac14\sum_i a_i^2\right).
\]
A partition state and its complement, with equal probabilities, have
all means $1/2$ and $S=A/2$ deterministically. Thus a YES instance has
lower envelope at most $b+\epsilon^2/8$. In a NO instance every binary
state has integer $S\ne A/2$, so $\operatorname{Var}(S)\ge1$ under
every feasible law. Its lower envelope is at least $b+\epsilon^2/2$.
The rational threshold $q=b+\epsilon^2/4$ separates the cases.
All numbers have polynomial encoding length, since $\log A$ is
polynomial in the original input length. For a fixed $\eta>0$, choose
fixed $K$ with $1/(16K)\le\eta$ to enforce the narrower interval.

For NP membership, the vertex-law LP has equality rows
$1,z_1,\ldots,z_n$ and right-hand side $(1,1/2,\ldots,1/2)$.
An optimal basic solution has at most $n+1$ positive probabilities.
It provides that many binary states and rational weights of polynomial
bit length: a nonsingular subsystem has zero-one entries and
half-integer right-hand side, so determinant bounds apply. Each vertex
product has polynomial bit length as a product of at most $n$ input
rationals. Exact arithmetic verifies nonnegative weights, total mass,
all means, and expected product at most $q$ in polynomial time.
\end{proof}

The NO case also has an explicit affine certificate in normalized
coordinates:
\begin{equation}\label{eq:partition-dual}
 \ell(z)=1-\epsilon^2A^2/8+\epsilon^2/2
    +\sum_i\left[\epsilon a_i+\frac{\epsilon^2}{2}(Aa_i-a_i^2)\right]z_i.
\end{equation}
Subtracting it from \eqref{eq:partition-expansion} gives
$\epsilon^2[(S-A/2)^2-1]/2+R(z)\ge0$ at every binary vertex of a
NO instance, hence everywhere by multiaffine interpolation. Its midpoint
value is $b+\epsilon^2/2$.
The reduction is from weakly NP-complete PARTITION; it establishes
ordinary, not strong, NP-hardness. It concerns varying unequal boxes
contained in a narrow common box, not the envelope on that one common
box. The separating accuracy can be exponentially small in input length.

\begin{corollary}[Graph-hull membership]\label{cor:single-product-membership}
Deciding whether $(\bar x,q)$ belongs to the scalar graph hull of the
same single monomial is NP-complete under the same narrow-box restrictions.
\end{corollary}
\begin{proof}
At the midpoint its common-threshold upper value is
$C=(1+\prod_i u_i)/2$. Retaining the product's nonnegative quadratic
terms gives
\[
 C\ge1+\epsilon A/2+\epsilon^2(A^2-\sum_i a_i^2)/4,
 \qquad C-b\ge\epsilon^2A^2/8\ge\epsilon^2/2>q-b.
\]
Thus the upper-hull condition always holds on the reduced instances,
and membership is exactly the lower-envelope threshold question.
Membership lies in NP: a basic distribution for the normalization,
coordinate, and value equations has at most $n+2$ positive weights.
These equations have rational polynomial-bit entries and determinant
bounds give polynomial-bit weights. Exact verification proves membership.
\end{proof}

\begin{proposition}[A fixed number of ratio types]\label{prop:ratio-types}
For one product on a rational strictly positive box having at most $k$
distinct coordinate aspect ratios, its exact scalar envelope and a
supporting affine minorant can be computed in $n^{O(k)}$ times a
polynomial in the input bit length. In particular fixed $k$ is tractable.
\end{proposition}
\begin{proof}
After substituting fixed coordinates and normalizing endpoints, a
binary product with count $h_j$ in ratio group $j$ has value
$(\prod_i\ell_i)\prod_{j=1}^k r_j^{h_j}$. To minimize this value minus
$\sum_i\pi_i z_i$ for arbitrary dual coefficients, sort the coefficients
within each group. For fixed counts choose the $h_j$ largest coefficients,
and use prefix sums. Enumerate all $\prod_j(n_j+1)\le(n+1)^k$ count
tuples. Rational multiplication and sorting have polynomial bit cost;
this is an exact vertex oracle. The determinant-bounded dual argument
of Corollary~\ref{cor:scalar-envelope-algorithm} then computes the lower
envelope and its minorant. The common-threshold upper envelope is
explicit. Thus mere closeness of the ratios does not substitute for a
bound on their number of distinct values.
\end{proof}

\subsection{Rank-one transport and logarithmic accuracy dependence}
Multimarginal optimal transport with specified finite marginals minimizes
expected cost over their joint laws. With $n$ fair binary marginals and
factored cost tensor
\[
 C=\bigotimes_{i=1}^n(1,1+\epsilon a_i),\qquad
 C_z=\prod_i(1+\epsilon a_i z_i),
\]
it is exactly the midpoint envelope above. This tensor is strictly
positive and rank one, with no sparse correction, and
\[
 1\le C_z\le\prod_i(1+\epsilon a_i)
       \le\sum_{j\ge0}(\epsilon A)^j\le16/15.
\]
\begin{corollary}[Rank-one value and weighted dual accuracy]
\label{cor:rank-one-precision}
Exact threshold evaluation for these rank-one MOT instances is
NP-complete. Unless $P=NP$, neither their values nor their weighted
vertex minima admit a deterministic additive-$\delta$ approximation
algorithm polynomial in rational input length and
$\log(C_{\max}/\delta)$. For randomized bounded-error algorithms the
corresponding consequence requires $NP\not\subseteq BPP$.
\end{corollary}
\begin{proof}
Take $\delta=\epsilon^2/32$ in the YES and NO intervals of
Theorem~\ref{thm:single-product-hardness}. Their approximation intervals
remain disjoint and $\log(1/\delta)$ is polynomial in the PARTITION
input length. Such an algorithm would decide PARTITION.
The weighted vertex oracle is hard directly: set
\[
 \pi_i=\epsilon a_i+\frac{\epsilon^2}{2}(Aa_i-a_i^2),\qquad
 d_0=1-\epsilon^2A^2/8.
\]
Equation~\eqref{eq:partition-expansion} gives
\[
 C_z-\pi^{\mathsf T}z=d_0+\frac{\epsilon^2}{2}(S(z)-A/2)^2+R(z).
\]
Its minimum is at most $d_0+\epsilon^2/8$ in a YES instance and at
least $d_0+\epsilon^2/2$ in a NO instance. The same approximation margin
and polynomial encoding bounds apply. Exact MOT decision has the
NP certificates already given. If an input convention requires bit
lengths polynomial in the number of marginals, append a number of dummy
binary factors $(1,1)$ linear in the original input bit length. Give
them fair marginals and zero dual potentials. They change neither
rank, cost range, value, nor oracle minimum, and meet that convention.
\end{proof}

The published structured-MOT paper of \citet[Section 7.2, after
Theorem 7.4]{AltschulerBoixAdsera2023} asks whether constant-rank
approximate dual minimization can have logarithmic inverse-accuracy
dependence, leading to exact algorithms. The preceding reduction rules
out that dependence in the rational bit model under the stated
complexity assumption, already at rank one. Their Corollary 7.5 has
polynomial dependence on inverse accuracy and is consistent with it.
This comparison identifies the historical question; it is not a claim
that no subsequent work addressed it. Nor is it a lower bound in an
unrestricted unit-cost arithmetic model. The unweighted smallest entry
of this tensor is simply one; arbitrary dual weights are essential.
None of these arguments proves hardness at fixed additive error.

\subsection{A boundary of multilinearity}\label{subsec:rational-powers}
\begin{example}[Two positive rational-power terms]\label{ex:rational-powers}
On the fixed interval $[1,3]$, let
$f_\epsilon(x)=x^{1-\epsilon}+x^{1+\epsilon}$, $0<\epsilon<1$.
The first term is concave and the second convex, while
\[
 f_\epsilon''(x)=\epsilon x^{-1-\epsilon}
       [(1+\epsilon)x^{2\epsilon}-(1-\epsilon)]>0.
\]
At $x=2$ their separate widths and full width are consequently
\[
 \tgap_\epsilon=3\sinh(\epsilon\log3)-4\sinh(\epsilon\log2),
 \qquad
 \hgap_\epsilon=1+3\cosh(\epsilon\log3)-4\cosh(\epsilon\log2).
\]
These follow by subtracting each function from its endpoint secant,
with the opposite subtraction for the concave term. Taylor expansion
gives
\[
 \tgap_\epsilon=(3\log3-4\log2)\epsilon+O(\epsilon^3),\qquad
 \hgap_\epsilon=\tfrac12[3(\log3)^2-4(\log2)^2]\epsilon^2
                                      +O(\epsilon^4).
\]
The first leading coefficient is $\log(27/16)>0$; the second is
positive by strict convexity of $x(\log x)^2$ on $[1,3]$.
Thus their ratio diverges like a positive constant times $1/\epsilon$.
Taking $\epsilon=1/k$, $k\ge2$, keeps rational exponents in
$[1/2,3/2]$ with only two positive unit coefficients and one variable.
Under $x=e^y$, however, both terms become convex exponentials on
$[0,\log3]$. Their lower envelopes are exact and endpoint secants
add, so the termwise and full widths coincide in that representation.
The example concerns the original coordinate and does not preclude
useful logarithmic reformulation.
\end{example}
